\documentclass{./worksheet}
\usepackage[colorlinks=true, urlcolor=blue]{hyperref}

\mytitle{Number Representations}

\begin{document}

This worksheet will give you some practice converting and working with decimal, binary, and hexadecimal numbers.
You can also raise your hand and ask for help if you are stuck!
\begin{enumerate}

\item Convert the decimal number $17$ to binary and hexadecimal.

\vspace{0.5in}

\item Convert the hexadecimal number $\texttt{0xf}$ to binary and decimal.

\vspace{0.5in}

\item Convert the binary number $\texttt{0b1010}$ to decimal and hexadecimal.

\vspace{0.5in}

\item Look at your answers to question 2 and 3. How many binary digits is needed to represent a single hexadecimal digit?

\item If a byte is 8 bits, how many hexadecimal digits does it take to represent a byte?

\item \textbf{Challenge}: Now that you know how to convert between decimal, binary, and hexadecimal, try applying bitwise operators onto these numbers:
\begin{itemize}
\item $\texttt{0b1 << 2} =$
\item $\texttt{8 \& 1} =$
\item $\texttt{0xa | 2} =$
\end{itemize}

\vspace{0.5in}

\item Adding and subtracting numbers in different number systems can be a bit
tricky at first, but it follows the same general rules as adding and subtracting in base 10.

Alternatively, you can convert back to base 10 and then perform the operation.

\textbf{Challenge}: Answer the following questions

\begin{itemize}
\item $\texttt{0b0010} + \texttt{0b1001} =$
\item $\texttt{0xa} - \texttt{0x5} =$
\item $\texttt{0xf} + 1 =$
\item $\texttt{0b1100} - \texttt{0xb} =$
\end{itemize}

\end{enumerate}


\end{document}
